\chapter{Use cases}
\label{ch:bijlage-usecases}

Deze bijlage bevat de volledige uitwerking van de use cases die werden opgesteld tijdens de analyse van de PoC in fase 3 (zie hoofdstuk~\ref{ch:fase3}).

\section{Receptiebon verwerken}
\label{uc:receptiebon-verwerken}

    \vspace{0.5em}

    \begin{itemize}[leftmargin=3.5cm, labelsep=0.5cm, nosep]
        \item[\textbf{Primaire actor:}] Productiechef, financieel bediende
        \item[\textbf{Stakeholders:}] CEO
        \item[\textbf{Precondities:}] De gebruiker is ingelogd
        \item[\textbf{Postcondities:}] Een receptiebon is verwerkt
    \end{itemize}

    \paragraph{Normaal verloop}
    \begin{enumerate}[nosep]
        \item De gebruiker wenst een receptiebon te verwerken.
        \item De gebruiker geeft een referentiecode in.
        \item De gebruiker kiest de juiste leverancier.
        \item De gebruiker kiest alle leeggoedtypes die van toepassing zijn.
        \item de gebruiker kiest de hoeveelheden voor elk leeggoedtype.
        \item Het systeem valideert volgens DR\_RECEPTIEBON.
        \item Het systeem registreert de receptiebon en de bijhorende stockbewegingen.
        \item Het systeem werkt het tegoed van de leverancier en de fysieke voorraad bij.
        \item Het systeem toont het detailscherm van de zonet aangemaakte receptiebon.
    \end{enumerate}

    \paragraph{Alternatieve verlopen}
    \begin{itemize}[nosep]
        \item[5A.] Het systeem detecteert dat de validatie volgens DR\_RECEPTIEBON faalt.
        \begin{itemize}
            \item[5A1.] Het systeem toont een gepaste melding.
            \item[5A2.] We keren terug naar stap 2.
        \end{itemize}
    \end{itemize}

    \paragraph{Domeinspecifieke regels}
    DR\_RECEPTIEBON
    \begin{itemize}
        \item referentiecode
        \begin{itemize}
            \item verplicht
            \item exact 6 tekens lang
        \end{itemize}
        \item leverancier
        \begin{itemize}
            \item verplicht
            \item moet bestaan in het systeem
        \end{itemize}
        \item leeggoedbewegingen
        \begin{itemize}
            \item verplicht
            \item aantal moet groter zijn dan nul
        \end{itemize}
    \end{itemize}

    \vspace{1cm}
\pagebreak
\section{Historiek bekijken}
\label{uc:historiek-bekijken}

    \vspace{0.5em}

    \begin{itemize}[leftmargin=3.5cm, labelsep=0.5cm, nosep]
        \item[\textbf{Primaire actor:}] Productiechef, financieel bediende, CEO
        \item[\textbf{Stakeholders:}] /
        \item[\textbf{Precondities:}] De gebruiker is ingelogd
        \item[\textbf{Postcondities:}] De gebruiker kan de historiek van stockbewegingen bezichtigen.
    \end{itemize}

    \paragraph{Normaal verloop}
    \begin{enumerate}[nosep]
        \item de gebruiker wil de historiek van stockbewegingen bekijken.
        \item De gebruiker navigeert naar het scherm voor historiek.
        \item De gebruiker kiest het gewenste leeggoedtype waar de historiek van bezichtigd wil worden.
        \item Het systeem toont alle historiek voor het gekozen leeggoedtype.
    \end{enumerate}

    \paragraph{Alternatieve verlopen}
    /

    \paragraph{Domeinspecifieke regels}
    /

    \vspace{1cm}

\section{Tegoed bekijken}
\label{uc:tegoed-bekijken}

    \vspace{0.5em}

    \begin{itemize}[leftmargin=3.5cm, labelsep=0.5cm, nosep]
        \item[\textbf{Primaire actor:}] Productiechef, financieel bediende, CEO
        \item[\textbf{Stakeholders:}] /
        \item[\textbf{Precondities:}] De gebruiker is ingelogd
        \item[\textbf{Postcondities:}] De gebruiker kan de stock van alle leeggoedtypes voor elke leverancier bezichtigen.
    \end{itemize}

    \paragraph{Normaal verloop}
    \begin{enumerate}[nosep]
        \item de gebruiker wil het tegoed van alle leeggoedtypes per leverancier bekijken.
        \item De gebruiker navigeert naar het scherm voor stockoverzicht.
        \item Het systeem toont de totale stock van elk leeggoedtype per leverancier.
    \end{enumerate}

    \paragraph{Alternatieve verlopen}
    /

    \paragraph{Domeinspecifieke regels}
    /

    \vspace{1cm}
\pagebreak
\section{Stock filteren}
\label{uc:stock-filteren}

    \vspace{0.5em}

    \begin{itemize}[leftmargin=3.5cm, labelsep=0.5cm, nosep]
        \item[\textbf{Primaire actor:}] Productiechef, financieel bediende, CEO
        \item[\textbf{Stakeholders:}] /
        \item[\textbf{Precondities:}] De gebruiker is ingelogd
        \item[\textbf{Postcondities:}] De gebruiker kan een gefilterde lijst van stock bezichtigen.
    \end{itemize}

    \paragraph{Normaal verloop}
    \begin{enumerate}[nosep]
        \item de gebruiker wil het tegoed voor een specifieke leverancier bezichtigen.
        \item De gebruiker navigeert naar het scherm voor stockbeheer.
        \item Het systeem toont de totale stock van elk leeggoedtype per leverancier.
        \item De gebruiker kiest de gewenste filters.
        \item Het systeem valideert de filters volgens DR\_STOCK\_FILTERS.
        \item het systeem toont alle stock voor de gekozen filters.
    \end{enumerate}

    \paragraph{Alternatieve verlopen}
    \begin{itemize}[nosep]
        \item[5A.] De validatie volgens DR\_STOCK\_FILTERS faalt.
        \begin{itemize}
            \item[5A1.] Het systeem toont een gepaste melding.
            \item[5A2.] We keren terug naar stap 4
        \end{itemize}

        \item[5B.] Het systeem detecteert dat voor de gekozen filters geen resultaten zijn.
        \begin{itemize}
            \item[5B1.] Het systeem toont een gepaste melding.
        \end{itemize}
    \end{itemize}

    \paragraph{Domeinspecifieke regels}
    DR\_STOCK\_FILTERS
    \begin{itemize}
        \item Leverancier
        \begin{itemize}
            \item gekozen leverancier moet bestaan in het systeem
        \end{itemize}
        \item Leeggoedtype
        \begin{itemize}
            \item gekozen leeggoedtype moet bestaan in het systeem
        \end{itemize}
    \end{itemize}

    \vspace{1cm}
\pagebreak
\section{Stockcorrectie uitvoeren}
\label{uc:stockcorrectie-uitvoeren}

    \begin{itemize}[leftmargin=3.5cm, labelsep=0.5cm, nosep]
        \item[\textbf{Primaire actor:}] Productiechef, financieel bediende, CEO
        \item[\textbf{Stakeholders:}] /
        \item[\textbf{Precondities:}] De gebruiker is ingelogd
        \item[\textbf{Postcondities:}] De gebruiker heeft de stock van een leeggoedtype voor een specifieke leverancier aangepast.
    \end{itemize}

    \paragraph{Normaal verloop}
    \begin{enumerate}[nosep]
        \item De gebruiker wil de stock van een leeggoedtype voor een specifieke leverancier aanpassen
        \item De gebruiker navigeert naar het scherm voor stockcorrecties.
        \item De gebruiker geeft het gewenste leeggoedtype en de leverancier in waarvoor een stockcorrectie uitgevoerd moet worden.
        \item Het systeem valideert volgens DR\_STOCKCORRECTIE\_FILTERS.
        \item De gebruiker geeft de gegevens van de correctie in.
        \item Het systeem valideert volgens DR\_STOCKCORRECTIE.
        \item Het systeem voert de stockcorrectie uit.
        \item Het systeem toont een gepaste melding.
    \end{enumerate}

    \paragraph{Alternatieve verlopen}
    \begin{itemize}[nosep]
        \item[6A.] het systeem detecteerd dat de validatie volgens DR\_STOCKCORRECTIE\_FILTERS faalt.
        \begin{itemize}
            \item[6A1.] Het systeem toont een gepaste melding.
            \item[6A2.] We keren terug naar stap 5
        \end{itemize}
        \item[6B.] Het systeem detecteert dat voor de gekozen filters geen resultaten zijn.
        \begin{itemize}
            \item[6B1.] Het systeem toont een gepaste melding.
            \item[6B2.] We keren terug naar stap 5
        \end{itemize}
        \item[8A.] De validatie volgens DR\_STOCKCORRECTIE faalt.
        \begin{itemize}
            \item[8A1.] Het systeem toont een gepaste melding.
            \item[8A2.] We keren terug naar stap 7
        \end{itemize}

        \item[8B.] Het systeem detecteert een verdachte correctie volgens DR\_STOCKCORRECTIE\_VERDACHT
        \begin{itemize}
            \item[8B1.] Het systeem vraagt om een bevesteging.
            \begin{itemize}
                \item[8B1A.] De gebruiker bevestigd de actie.
                \begin{itemize}
                    \item[8B1A1.] Het systeem toont een gepaste melding
                    \item[8B1A2.] We gaan verder naar stap 9
                \end{itemize}
                \item[8B1B.] De gebruiker bevestigd niet.
                \begin{itemize}
                    \item[8B1B1.] Het systeem toont een gepaste melding.
                    \item[8B1B2.] We keren terug naar stap 7
                \end{itemize}
            \end{itemize}
        \end{itemize}
    \end{itemize}

    \paragraph{Domeinspecifieke regels}
    DR\_STOCKCORRECTIE\_FILTERS
    \begin{itemize}
        \item Leverancier
        \begin{itemize}
            \item gekozen leverancier moet bestaan in het systeem
        \end{itemize}
        \item Leeggoedtype
        \begin{itemize}
            \item gekozen leeggoedtype moet bestaan in het systeem
        \end{itemize}
    \end{itemize}
    \\
    DR\_STOCKCORRECTIE
     \begin{itemize}
        \item Hoeveelheid
        \begin{itemize}
            \item verplicht
            \item moet een geheel getal zijn
            \item mag niet nul zijn
        \end{itemize}
    \end{itemize}
    \\

    \vspace{1cm}

\section{Leverancier aanmaken}
\label{uc:leverancier-aanmaken}

\vspace{0.5em}

\begin{itemize}[leftmargin=3.5cm, labelsep=0.5cm, nosep]
    \item[\textbf{Primaire actor:}] Financieel bediende, Productiechef, CEO
    \item[\textbf{Stakeholders:}] /
    \item[\textbf{Precondities:}] De gebruiker is ingelogd
    \item[\textbf{Postcondities:}] Een nieuwe leverancier is toegevoegd aan het systeem.
\end{itemize}

\paragraph{Normaal verloop}
\begin{enumerate}[nosep]
    \item De gebruiker wil een nieuwe leverancier aanmaken.
    \item De gebruiker navigeert naar het scherm voor leveranciers.
    \item De gebruiker kiest de optie om een nieuwe leverancier toe te voegen.
    \item De gebruiker vult de naam en het adres van de leverancier in.
    \item Het systeem valideert volgens DR\_LEVERANCIER.
    \item Het systeem registreert de leverancier.
\end{enumerate}

\paragraph{Alternatieve verlopen}
\begin{itemize}[nosep]
    \item[5A.] De validatie volgens DR\_LEVERANCIER faalt.
    \begin{itemize}
        \item[5A1.] Het systeem toont een gepaste melding.
        \item[5A2.] We keren terug naar stap 4.
    \end{itemize}
\end{itemize}

\paragraph{Domeinspecifieke regels}
DR\_LEVERANCIER
\begin{itemize}
    \item naam
    \begin{itemize}
        \item verplicht
    \end{itemize}
    \item adres
    \begin{itemize}
        \item verplicht
    \end{itemize}
\end{itemize}

\vspace{1cm}
\pagebreak

\section{Leeggoedtype aanmaken}
\label{uc:leeggoedtype-aanmaken}

\vspace{0.5em}

\begin{itemize}[leftmargin=3.5cm, labelsep=0.5cm, nosep]
    \item[\textbf{Primaire actor:}] Financieel bediende, Productiechef, CEO
    \item[\textbf{Stakeholders:}] /
    \item[\textbf{Precondities:}] De gebruiker is ingelogd
    \item[\textbf{Postcondities:}] Een nieuw leeggoedtype is toegevoegd aan het systeem.
\end{itemize}

\paragraph{Normaal verloop}
\begin{enumerate}[nosep]
    \item De gebruiker wil een nieuw leeggoedtype aanmaken.
    \item De gebruiker navigeert naar het scherm voor leeggoedtypes.
    \item De gebruiker kiest de optie om een nieuw leeggoedtype toe te voegen.
    \item De gebruiker geeft de artikelcode en de naam van het leeggoedtype in.
    \item Het systeem valideert volgens DR\_LEEGGOEDTYPE.
    \item Het systeem registreert het leeggoedtype en toont een gepaste melding.
\end{enumerate}

\paragraph{Alternatieve verlopen}
\begin{itemize}[nosep]
    \item[5A.] De validatie volgens DR\_LEEGGOEDTYPE faalt.
    \begin{itemize}
        \item[5A1.] Het systeem toont een gepaste melding.
        \item[5A2.] We keren terug naar stap 4.
    \end{itemize}
\end{itemize}

\paragraph{Domeinspecifieke regels}
DR\_LEEGGOEDTYPE
\begin{itemize}
    \item artikelcode
    \begin{itemize}
        \item verplicht
        \item moet voldoen aan het regex patroon \texttt{LG\textbackslash d\{3\}}
        \item moet uniek zijn binnen het systeem
    \end{itemize}
    \item naam
    \begin{itemize}
        \item verplicht
        \item moet uniek zijn binnen het systeem
    \end{itemize}
\end{itemize}

\vspace{1cm}

\section{Inloggen}
\label{uc:inloggen}

    \vspace{0.5em}

    \begin{itemize}[leftmargin=3.5cm, labelsep=0.5cm, nosep]
        \item[\textbf{Primaire actor:}] Productiechef, financieel bediende, CEO
        \item[\textbf{Stakeholders:}] /
        \item[\textbf{Precondities:}] /
        \item[\textbf{Postcondities:}] De gebruiker is aangemeld
    \end{itemize}

    \paragraph{Normaal verloop}
    \begin{enumerate}[nosep]
        \item De gebruiker wenst zich aan te melden.
        \item Het systeem vraagt email en wachtwoord.
        \item De gebruiker voert de gegevens in.
        \item Het systeem valideert volgens DR\_AANMELDEN
    \end{enumerate}

    \paragraph{Alternatieve verlopen}
    \begin{itemize}[nosep]
        \item[4A.] Het systeem detecteert dat de email en/of wachtwoord niet ingevuld zijn
        \begin{itemize}
            \item[4A1.] Het systeem toont een gepaste melding.
            \item[4A2.] We keren terug naar stap 2
        \end{itemize}

        \item[4B.] Het systeem detecteert dat de email niet overeenstemt met het bijhorend wachtwoord
        \begin{itemize}
            \item[4B1.] Het systeem toont een gepaste melding.
            \item[4B1.] We keren terug naar stap 2
        \end{itemize}
    \end{itemize}

    \paragraph{Domeinspecifieke regels}
    DR\_AANMELDEN
    \begin{itemize}
        \item email en wachtwoord zijn verplicht
        \item email en wachtwoord moeten overeenstemmen.
    \end{itemize}
